\section{From a structural ceiling to finite-time access}\label{sec:introduction}
The reciprocal chiral exchange candidate couples a native state to a periodic carrier and a reversible receiver through two finite distributed profiles. Its original reciprocal feedback provides well-posed evolution and matching signed storage accounts~\cite{p1}. A subsequent theorem establishes later native credit on a specified initially empty solution~\cite{p2}. The projection-capacity analysis then supplies an additional native invariant and a sharper structural ceiling, while a separate sufficient regime yields high receiving and recovery fractions relative to an intermediate transferred stock~\cite{p3}. These results distinguish how much can leave from what fraction of a particular transferred amount can return.

The remaining gap addressed here is quantitative access: does an actual unchanged solution reach a definite fraction of the joint ceiling at a finite time? A sharp minimum on an invariant level need not be reached dynamically, and high recovered fractions can coexist with small absolute transfer. For the fixed-yaw coherent empty-launch family, we establish access to more than one quarter of the joint ceiling under an explicit sufficient condition. Geometry and all coefficients are fixed before choosing the phase. The time and the fractional guarantee are uniform over the declared open phase range.

The accepted investigation~\cite{fa} and its project-internal scientific review~\cite{fareview} are the sources of the new theorem and full proof. Section~\ref{sec:model} fixes the model and inherited ceiling. Section~\ref{sec:theorem} states the certificate; Section~\ref{sec:proof} proves it through an original-variable energy estimate, a drive-dependent Volterra bound and a reference residual. Section~\ref{sec:regime} proves nonemptiness of the finite parameter regime. Section~\ref{sec:interpretation} separates the resulting stock guarantee from receiving, recovery and geometric questions. Appendix~\ref{sec:units} checks dimensions without assigning a physical clock.

The motivating picture retains a central starting or meeting region, opposed cones, outward spreading and a curved around-and-inward return, with stationary supports and possible transfer elsewhere. The present theorem is a storage statement for the specified candidate. A norm estimate does not assign this intended spatial route.
